\documentclass[11pt,a4paper]{article}

% Packages
\usepackage[utf8]{inputenc}
\usepackage[T1]{fontenc}
\usepackage{amsmath,amssymb,amsthm}
\usepackage{graphicx}
\usepackage{booktabs}
\usepackage{multirow}
\usepackage{geometry}
\usepackage{hyperref}
\usepackage{enumitem}
\usepackage{array}
\usepackage{longtable}

% Page geometry
\geometry{margin=2.5cm}

% Title information
\title{\textbf{REC-Level Battery Optimization with Heterogeneous Distributed Storage}}
\author{REC Optimization Team}
\date{February 1, 2026\\
\vspace{0.3cm}
\textit{Document Version: 1.0}\\
\textit{Scenario: C3 - Single Supplier with REC and Distributed Battery Optimization}\\
\textit{Network: 1-LV-rural3--2-no\_sw (9 nodes, 3 prosumers with PV+Storage)}}

\begin{document}

\maketitle

\section{Executive Summary}

This document describes the Mixed-Integer Linear Programming (MILP) optimization model for coordinating \textbf{three heterogeneous battery storage systems} distributed across prosumer nodes in a Renewable Energy Community (REC). The optimization is implemented in \textbf{Pyomo} and solved with \textbf{Gurobi}. The optimizer operates at the REC level to minimize total community energy costs while respecting:

\begin{itemize}
    \item Individual battery technical specifications (capacity, power, efficiency)
    \item Physical distribution of assets (batteries remain at prosumer locations)
    \item Energy balance constraints at both node and REC levels
    \item Grid interaction limits and market price signals
\end{itemize}

\textbf{Key Innovation:} Centralized REC-level optimization of heterogeneous distributed batteries (not a single centralized battery), executed using \textbf{intraday (ID) market forecasts only} with rolling horizon updates every hour.

\section{System Architecture}

\subsection{REC Composition}

\begin{table}[h]
\centering
\small
\begin{tabular}{@{}cccccl@{}}
\toprule
\textbf{Node} & \textbf{Type} & \textbf{Load Profile} & \textbf{PV Profile} & \textbf{Storage Profile} & \textbf{Battery} \\ 
\midrule
1 & Consumer & G4 & -- & -- & No \\
2 & Prosumer & G6 & PV4 & Storage\_PV4\_H0-G & \textbf{Yes} \\
3 & Consumer & H0-L & -- & -- & No \\
4 & Consumer & H0-L & -- & -- & No \\
5 & Consumer & H0-G & -- & -- & No \\
6 & Prosumer & H0-L & PV3 & Storage\_PV3\_H0-L & \textbf{Yes} \\
7 & Consumer & G1 & -- & -- & No \\
8 & Prosumer & H0-L & PV1 & Storage\_PV1\_H0-L & \textbf{Yes} \\
9 & Consumer & H0-L & -- & -- & No \\
\bottomrule
\end{tabular}
\caption{REC participant composition}
\end{table}

\textbf{Total REC Capacity:}
\begin{itemize}
    \item \textbf{Consumers:} 6 nodes (42,832 kWh/year)
    \item \textbf{Prosumers:} 3 nodes (21,069 kWh/year load, 25,083 kWh/year PV)
    \item \textbf{Net Annual Balance:} $\sim$3,900 kWh surplus (REC is net exporter)
\end{itemize}

\section{Heterogeneous Battery Specifications}

\subsection{Node 2: Fire Fighting Station (Commercial)}

\textbf{Rationale:} Large commercial PV system (17.68 kWp) requires substantial storage for load shifting and grid arbitrage.

\begin{table}[h]
\centering
\small
\begin{tabular}{@{}llll@{}}
\toprule
\textbf{Parameter} & \textbf{Value} & \textbf{Unit} & \textbf{Notes} \\ 
\midrule
Capacity & 40 & kWh & 2.3h at peak PV output \\
Max Charge Power & 20 & kW & $1.13\times$ peak PV (17.68 kWp) \\
Max Discharge Power & 20 & kW & Can supply full load + export \\
Charge Efficiency & 95 & \% & Commercial-grade Li-ion \\
Discharge Efficiency & 95 & \% & Round-trip: 90.25\% \\
Self-Discharge Rate & 0.1 & \%/hour & 2.4\% per day \\
SOC Min & 10 & \% & Deep discharge capable \\
SOC Max & 100 & \% & Full capacity utilization \\
Initial SOC & 50 & \% & Starting condition \\
Usable Capacity & 36 & kWh & $(100\% - 10\%) \times 40$ kWh \\
\bottomrule
\end{tabular}
\caption{Node 2 battery specifications}
\end{table}

\textbf{Strategic Role:} Primary grid arbitrage battery, handles bulk energy shifting, peak shaving for REC.

\subsection{Node 6: Household (Medium Residential)}

\textbf{Rationale:} Medium PV system (4.2 kWp) serving high residential load (14,094 kWh/year).

\begin{table}[h]
\centering
\small
\begin{tabular}{@{}llll@{}}
\toprule
\textbf{Parameter} & \textbf{Value} & \textbf{Unit} & \textbf{Notes} \\ 
\midrule
Capacity & 10 & kWh & 2.4h at peak PV output \\
Max Charge Power & 5 & kW & $1.19\times$ peak PV (4.2 kWp) \\
Max Discharge Power & 5 & kW & Standard residential inverter \\
Charge Efficiency & 92 & \% & Mid-tier residential battery \\
Discharge Efficiency & 92 & \% & Round-trip: 84.64\% \\
Self-Discharge Rate & 0.2 & \%/hour & 4.8\% per day \\
SOC Min & 20 & \% & Battery protection threshold \\
SOC Max & 100 & \% & Full capacity utilization \\
Initial SOC & 50 & \% & Starting condition \\
Usable Capacity & 8 & kWh & $(100\% - 20\%) \times 10$ kWh \\
\bottomrule
\end{tabular}
\caption{Node 6 battery specifications}
\end{table}

\textbf{Strategic Role:} Evening residential peak support, local self-consumption optimization.

\subsection{Node 8: Household (Small Residential)}

\textbf{Rationale:} Small PV system (2.6 kWp) serving low residential load (1,804 kWh/year).

\begin{table}[h]
\centering
\small
\begin{tabular}{@{}llll@{}}
\toprule
\textbf{Parameter} & \textbf{Value} & \textbf{Unit} & \textbf{Notes} \\ 
\midrule
Capacity & 6.5 & kWh & 2.5h at peak PV output \\
Max Charge Power & 3.25 & kW & $1.25\times$ peak PV (2.6 kWp) \\
Max Discharge Power & 3.25 & kW & Basic residential inverter \\
Charge Efficiency & 90 & \% & Entry-level battery \\
Discharge Efficiency & 90 & \% & Round-trip: 81\% \\
Self-Discharge Rate & 0.3 & \%/hour & 7.2\% per day \\
SOC Min & 20 & \% & Battery protection threshold \\
SOC Max & 100 & \% & Full capacity utilization \\
Initial SOC & 50 & \% & Starting condition \\
Usable Capacity & 5.2 & kWh & $(100\% - 20\%) \times 6.5$ kWh \\
\bottomrule
\end{tabular}
\caption{Node 8 battery specifications}
\end{table}

\textbf{Strategic Role:} Local self-consumption only, limited REC-level contribution due to low efficiency.

\subsection{Aggregate Battery Fleet}

\begin{table}[h]
\centering
\begin{tabular}{@{}lll@{}}
\toprule
\textbf{Metric} & \textbf{Total} & \textbf{Notes} \\ 
\midrule
Total Nominal Capacity & 56.5 kWh & Sum of all three batteries \\
Total Usable Capacity & 49.2 kWh & Accounting for SOC limits \\
Total Charge Power & 28.25 kW & Sum of max charge rates \\
Total Discharge Power & 28.25 kW & Sum of max discharge rates \\
Weighted Avg Efficiency & 87.8\% & Capacity-weighted round-trip \\
\bottomrule
\end{tabular}
\caption{Aggregate battery fleet specifications}
\end{table}

\section{Mathematical Model Formulation}

\subsection{Sets and Indices}

\begin{itemize}
    \item $\mathcal{N} = \{1, 2, 3, 4, 5, 6, 7, 8, 9\}$: All REC nodes
    \item $\mathcal{B} = \{2, 6, 8\} \subset \mathcal{N}$: Nodes with batteries (prosumers)
    \item $\mathcal{C} = \{1, 3, 4, 5, 7, 9\} \subset \mathcal{N}$: Nodes without batteries (consumers)
    \item $\mathcal{T} = \{0, 1, \ldots, 95\}$: Time intervals (15-min resolution, 24 hours)
\end{itemize}

\subsection{Parameters}

\subsubsection*{Time Series Data}
\begin{itemize}
    \item $L_{i,t}$: Load demand at node $i$, time $t$ [kW]
    \item $\text{PV}_{i,t}$: PV generation at node $i$, time $t$ [kW] ($i \in \mathcal{B}$ only)
    \item $\pi_t^{\text{ID}}$: Intraday market price at time $t$ [€/kWh]
    \item $\pi_t^{\text{FI}}$: Feed-in tariff at time $t$ [€/kWh]
\end{itemize}

\subsubsection*{Battery Specifications (for $i \in \mathcal{B}$)}
\begin{itemize}
    \item $E_i^{\text{cap}}$: Battery capacity [kWh]
    \item $P_i^{\text{ch,max}}$: Max charging power [kW]
    \item $P_i^{\text{dch,max}}$: Max discharging power [kW]
    \item $\eta_i^{\text{ch}}$: Charging efficiency [dimensionless]
    \item $\eta_i^{\text{dch}}$: Discharging efficiency [dimensionless]
    \item $\sigma_i$: Self-discharge rate [1/hour]
    \item $\text{SOC}_i^{\min}$: Minimum state of charge [\% of $E_i^{\text{cap}}$]
    \item $\text{SOC}_i^{\max}$: Maximum state of charge [\% of $E_i^{\text{cap}}$]
    \item $\text{SOC}_i^{\text{init}}$: Initial state of charge [\% of $E_i^{\text{cap}}$]
    \item $\text{SOC}_{i,0}$: Current state of charge at the start of each optimization stage [kWh]
\end{itemize}

\subsubsection*{Economic Parameters}
\begin{itemize}
    \item $\pi^{\text{grid}}$: Grid fee [€/kWh] = 0.02
    \item $\Delta t$: Time interval duration [hours] = 0.25 (15 minutes)
\end{itemize}

\subsection{Decision Variables}

\subsubsection*{Power Flows (Continuous, $\geq 0$)}
\begin{itemize}
    \item $P_{i,t}^{\text{grid,import}}$: Power imported from grid at node $i$, time $t$ [kW]
    \item $P_{i,t}^{\text{grid,export}}$: Power exported to grid at node $i$, time $t$ [kW]
    \item $P_{i,t}^{\text{rec,import}}$: Power imported from REC at node $i$, time $t$ [kW]
    \item $P_{i,t}^{\text{rec,export}}$: Power exported to REC at node $i$, time $t$ [kW]
\end{itemize}

\subsubsection*{Battery Variables ($i \in \mathcal{B}$ only)}
\begin{itemize}
    \item $P_{i,t}^{\text{ch}}$: Charging power [kW] (Continuous, $\geq 0$)
    \item $P_{i,t}^{\text{dch}}$: Discharging power [kW] (Continuous, $\geq 0$)
    \item $E_{i,t}^{\text{SOC}}$: State of charge [kWh] (Continuous, $\geq 0$)
\end{itemize}

\subsubsection*{Binary Variables}
\begin{itemize}
    \item $b_{i,t}^{\text{ch}}$: Battery charging indicator (1 = charging, 0 = not)
    \item $b_{i,t}^{\text{grid}}$: Grid flow direction (1 = import, 0 = export)
\end{itemize}

\subsection{Objective Function}

Minimize Total REC Cost using \textbf{intraday (ID) market prices only}. Battery optimization is performed with rolling horizon (1-hour ahead) using accurate ID forecasts, updated every hour.

\begin{equation}
\begin{split}
\min Z = \sum_{i \in \mathcal{N}} \sum_{t \in \mathcal{T}} \Big[
    &P_{i,t}^{\text{grid,import}} \cdot \pi_t^{\text{ID}} \cdot \Delta t \\
    &+ (P_{i,t}^{\text{grid,import}} + P_{i,t}^{\text{grid,export}}) \cdot \pi^{\text{grid}} \cdot \Delta t \\
    &- P_{i,t}^{\text{grid,export}} \cdot \pi_t^{\text{FI}} \cdot \Delta t
\Big]
\end{split}
\end{equation}

\textbf{Note:} Internal REC transactions ($P^{\text{rec,import}}$ and $P^{\text{rec,export}}$) cancel out in total cost calculation since one member's export is another's import at the same price (€0.08/kWh).

\textbf{Market Strategy:} Day-ahead (DA) market commitments are made \textit{without} battery optimization (baseline). Battery optimization occurs \textit{only} in the ID market where forecasts are more accurate.

\subsection{Constraints}

\subsubsection{Energy Balance -- Consumers ($i \in \mathcal{C}$)}

\begin{equation}
L_{i,t} = P_{i,t}^{\text{grid,import}} + P_{i,t}^{\text{rec,import}} - P_{i,t}^{\text{grid,export}} - P_{i,t}^{\text{rec,export}} \quad \forall i \in \mathcal{C}, \forall t \in \mathcal{T}
\end{equation}

Consumers have no local generation or storage, so load must equal net imports.

\subsubsection{Energy Balance -- Prosumers ($i \in \mathcal{B}$)}

\begin{equation}
\begin{split}
\text{PV}_{i,t} + P_{i,t}^{\text{dch}} + P_{i,t}^{\text{grid,import}} + P_{i,t}^{\text{rec,import}} \\
- L_{i,t} - P_{i,t}^{\text{ch}} - P_{i,t}^{\text{grid,export}} - P_{i,t}^{\text{rec,export}} = 0 \quad \forall i \in \mathcal{B}, \forall t \in \mathcal{T}
\end{split}
\end{equation}

Prosumers have \textbf{load consumption} ($L_{i,t}$), \textbf{PV generation}, and \textbf{battery storage}. Energy balance: all sources (PV + battery discharge + imports) must equal all sinks (load + battery charge + exports).

\subsubsection{REC Energy Balance (Community-Wide)}

\begin{equation}
\sum_{i \in \mathcal{N}} P_{i,t}^{\text{rec,export}} = \sum_{i \in \mathcal{N}} P_{i,t}^{\text{rec,import}} \quad \forall t \in \mathcal{T}
\end{equation}

Total energy exported within REC must equal total energy imported (conservation of energy within community).

\subsubsection{Battery State of Charge Dynamics ($i \in \mathcal{B}$)}

\textbf{Initial Condition ($t = 0$):}
\begin{equation}
E_{i,0}^{\text{SOC}} = \text{SOC}_{i,0} \quad \forall i \in \mathcal{B}
\end{equation}

\textbf{Evolution ($t > 0$):}
\begin{equation}
\begin{split}
E_{i,t}^{\text{SOC}} = &E_{i,t-1}^{\text{SOC}} \cdot (1 - \sigma_i \cdot \Delta t) \\
                       &+ P_{i,t-1}^{\text{ch}} \cdot \eta_i^{\text{ch}} \cdot \Delta t \\
                       &- P_{i,t-1}^{\text{dch}} / \eta_i^{\text{dch}} \cdot \Delta t 
\quad \forall i \in \mathcal{B}, \forall t \in \mathcal{T} \setminus \{0\}
\end{split}
\end{equation}

\textbf{Physical Interpretation:}
\begin{itemize}
    \item Battery loses energy over time due to self-discharge ($\sigma$)
    \item Charging adds energy with efficiency loss ($\eta^{\text{ch}} < 1$)
    \item Discharging removes more energy than delivered ($1/\eta^{\text{dch}} > 1$)
\end{itemize}

\subsubsection{Battery State of Charge Limits ($i \in \mathcal{B}$)}

\begin{equation}
E_i^{\text{cap}} \cdot \text{SOC}_i^{\min} \leq E_{i,t}^{\text{SOC}} \leq E_i^{\text{cap}} \cdot \text{SOC}_i^{\max} \quad \forall i \in \mathcal{B}, \forall t \in \mathcal{T}
\end{equation}

\textbf{Node-Specific Limits:}
\begin{itemize}
    \item Node 2: $4 \text{ kWh} \leq E_{2,t}^{\text{SOC}} \leq 40 \text{ kWh}$ (10\% -- 100\%)
    \item Node 6: $2 \text{ kWh} \leq E_{6,t}^{\text{SOC}} \leq 10 \text{ kWh}$ (20\% -- 100\%)
    \item Node 8: $1.3 \text{ kWh} \leq E_{8,t}^{\text{SOC}} \leq 6.5 \text{ kWh}$ (20\% -- 100\%)
\end{itemize}

\subsubsection{Battery Charging Power Limits ($i \in \mathcal{B}$)}

\begin{equation}
0 \leq P_{i,t}^{\text{ch}} \leq \min(P_i^{\text{ch,max}}, E_i^{\text{cap}} \cdot 1.0) \quad \forall i \in \mathcal{B}, \forall t \in \mathcal{T}
\end{equation}

\textbf{Node-Specific Limits:}
\begin{itemize}
    \item Node 2: $0 \leq P_{2,t}^{\text{ch}} \leq 20$ kW
    \item Node 6: $0 \leq P_{6,t}^{\text{ch}} \leq 5$ kW
    \item Node 8: $0 \leq P_{8,t}^{\text{ch}} \leq 3.25$ kW
\end{itemize}

\subsubsection{Battery Discharging Power Limits ($i \in \mathcal{B}$)}

\begin{equation}
0 \leq P_{i,t}^{\text{dch}} \leq \min(P_i^{\text{dch,max}}, E_i^{\text{cap}} \cdot 1.0) \quad \forall i \in \mathcal{B}, \forall t \in \mathcal{T}
\end{equation}

\textbf{Node-Specific Limits:}
\begin{itemize}
    \item Node 2: $0 \leq P_{2,t}^{\text{dch}} \leq 20$ kW
    \item Node 6: $0 \leq P_{6,t}^{\text{dch}} \leq 5$ kW
    \item Node 8: $0 \leq P_{8,t}^{\text{dch}} \leq 3.25$ kW
\end{itemize}

\subsubsection{No Simultaneous Charge/Discharge ($i \in \mathcal{B}$)}

\begin{align}
P_{i,t}^{\text{ch}} &\leq P_i^{\text{ch,max}} \cdot b_{i,t}^{\text{ch}} \quad \forall i \in \mathcal{B}, \forall t \in \mathcal{T} \\
P_{i,t}^{\text{dch}} &\leq P_i^{\text{dch,max}} \cdot (1 - b_{i,t}^{\text{ch}}) \quad \forall i \in \mathcal{B}, \forall t \in \mathcal{T} \\
b_{i,t}^{\text{ch}} &\in \{0, 1\} \quad \forall i \in \mathcal{B}, \forall t \in \mathcal{T}
\end{align}

\textbf{Physical Interpretation:} Battery cannot charge and discharge simultaneously. Binary variable $b_{i,t}^{\text{ch}}$ enforces this logical constraint.

\subsubsection{No Simultaneous Grid Import/Export ($i \in \mathcal{N}$)}

\begin{align}
P_{i,t}^{\text{grid,import}} &\leq M \cdot b_{i,t}^{\text{grid}} \quad \forall i \in \mathcal{N}, \forall t \in \mathcal{T} \\
P_{i,t}^{\text{grid,export}} &\leq M \cdot (1 - b_{i,t}^{\text{grid}}) \quad \forall i \in \mathcal{N}, \forall t \in \mathcal{T} \\
b_{i,t}^{\text{grid}} &\in \{0, 1\} \quad \forall i \in \mathcal{N}, \forall t \in \mathcal{T}
\end{align}

\textbf{Physical Interpretation:} Node cannot import from and export to grid simultaneously. $M$ is a sufficiently large constant (e.g., 1000 kW).

\section{Optimization Strategy and Expected Behavior}

\subsection{Rolling Horizon Optimization (Intraday Only)}

Battery optimization is performed \textbf{only in the intraday (ID) market} using a rolling horizon approach:

\begin{enumerate}
    \item \textbf{Day-Ahead (DA) Market:} REC makes day-ahead commitments \textit{without} battery optimization, based on DA load/generation forecasts.
    \item \textbf{Intraday (ID) Optimization:} Battery schedule is optimized using accurate ID forecasts (1-hour ahead horizon), refreshed every hour.
    \item \textbf{State Continuity:} The \textbf{current battery SOC} ($\text{SOC}_{i,0}$) from the previous executed schedule is used as the initial condition for each new optimization.
\end{enumerate}

\textbf{Rationale:} ID forecasts have significantly lower error than DA forecasts (RMSE reduced by $\sim$70\%), making battery optimization more effective and reducing imbalance costs.

This ensures feasibility and continuity by enforcing the actual battery state at the start of each rolling horizon optimization.

\subsection{Charging Priority (Morning/Midday)}

When REC has PV surplus, the optimizer charges batteries in this order:

\begin{enumerate}
    \item \textbf{Node 2 (40 kWh, 90.25\% efficiency):}
        \begin{itemize}
            \item Highest capacity $\rightarrow$ can absorb most surplus
            \item Best efficiency $\rightarrow$ minimizes energy loss
            \item Lowest SOC limit (10\%) $\rightarrow$ more usable capacity
            \item \textbf{Target fill time:} 2 hours at 20 kW = full charge
        \end{itemize}
    
    \item \textbf{Node 6 (10 kWh, 84.64\% efficiency):}
        \begin{itemize}
            \item Medium capacity $\rightarrow$ residential load matching
            \item Good efficiency $\rightarrow$ acceptable losses
            \item \textbf{Target fill time:} 2 hours at 5 kW = full charge
        \end{itemize}
    
    \item \textbf{Node 8 (6.5 kWh, 81\% efficiency):}
        \begin{itemize}
            \item Small capacity $\rightarrow$ limited absorption
            \item Lower efficiency $\rightarrow$ use sparingly
            \item \textbf{Target fill time:} 2 hours at 3.25 kW = full charge
        \end{itemize}
\end{enumerate}

\textbf{Rationale:} Fill largest and most efficient batteries first to maximize REC-level energy arbitrage opportunities.

\subsection{Discharging Priority (Evening Peak)}

When REC has load deficit (evening), the optimizer discharges batteries in this order:

\begin{enumerate}
    \item \textbf{Node 8 (small, low efficiency):}
        \begin{itemize}
            \item Discharge early $\rightarrow$ preserve Node 2 for bigger arbitrage
            \item Limited capacity $\rightarrow$ won't sustain long discharge
            \item \textbf{Target:} Local self-consumption only
        \end{itemize}
    
    \item \textbf{Node 6 (medium, good efficiency):}
        \begin{itemize}
            \item Support residential evening peak (18:00--21:00)
            \item Moderate capacity $\rightarrow$ 2--3 hours discharge at 5 kW
            \item \textbf{Target:} Residential load support
        \end{itemize}
    
    \item \textbf{Node 2 (large, high efficiency):}
        \begin{itemize}
            \item Reserve for highest-value discharge periods
            \item Large capacity $\rightarrow$ 4+ hours discharge at 20 kW
            \item \textbf{Target:} Grid arbitrage, peak demand reduction
        \end{itemize}
\end{enumerate}

\textbf{Rationale:} Preserve high-capacity, high-efficiency battery (Node 2) for maximum-value discharge opportunities (highest grid prices or peak demand charges).

\subsection{Daily Cycle Example}

\textbf{Typical Summer Day (High PV):}

\begin{table}[h]
\centering
\tiny
\begin{tabular}{@{}p{2.5cm}p{2.5cm}p{2cm}p{2cm}p{2cm}p{2cm}@{}}
\toprule
\textbf{Time} & \textbf{REC Net} & \textbf{Node 2} & \textbf{Node 6} & \textbf{Node 8} & \textbf{Grid} \\ 
\midrule
00:00--06:00 & Deficit (consumers) & Idle (preserve) & Idle & Idle & Import ($\sim$15 kW) \\
06:00--09:00 & Small surplus (PV) & \textbf{Charging} (0$\to$20 kW) & \textbf{Charging} (0$\to$5 kW) & \textbf{Charging} (0$\to$3 kW) & Minimal \\
09:00--12:00 & Large surplus (peak PV) & \textbf{Charging} (full by 11:00) & \textbf{Charging} (full by 11:00) & \textbf{Charging} (full by 11:00) & Export surplus \\
12:00--15:00 & Max surplus & Idle (full) & Idle (full) & Idle (full) & Export ($\sim$25 kW) \\
15:00--17:00 & Declining surplus & Idle & Idle & Idle & Minimal \\
17:00--19:00 & Deficit (no PV) & Idle (reserve) & \textbf{Discharging} (5 kW) & \textbf{Discharging} (3 kW) & Import ($\sim$7 kW) \\
19:00--21:00 & High deficit (peak) & \textbf{Discharging} (20 kW) & \textbf{Discharging} (5 kW) & Empty & Import ($\sim$10 kW) \\
21:00--24:00 & Moderate deficit & \textbf{Discharging} (10 kW) & Empty & Empty & Import ($\sim$5 kW) \\
\bottomrule
\end{tabular}
\caption{Daily battery operation cycle}
\end{table}

\textbf{Key Metrics:}
\begin{itemize}
    \item Total Battery Charged: $\sim$56 kWh (accounting for efficiency losses)
    \item Total Battery Discharged: $\sim$50 kWh delivered to loads
    \item Round-trip Energy Loss: $\sim$6 kWh (10--15\% depending on battery mix)
    \item Peak Grid Import Reduction: $\sim$25 kW (batteries discharge during peak)
    \item Peak Grid Export Reduction: $\sim$15 kW (batteries charge during surplus)
\end{itemize}

\section{Expected Outcomes}

\textbf{Baseline (No Batteries):}
\begin{itemize}
    \item Total daily cost: $\sim$€X (grid purchases -- feed-in revenue)
    \item Peak grid import: $\sim$40 kW
    \item Peak grid export: $\sim$30 kW
\end{itemize}

\textbf{With Optimized Batteries:}
\begin{itemize}
    \item Total daily cost: $\sim$€Y (expect 15--25\% reduction)
    \item Peak grid import: $\sim$20 kW (50\% reduction via discharge)
    \item Peak grid export: $\sim$20 kW (33\% reduction via charge)
    \item Self-consumption ratio: +20--30\% (PV $\rightarrow$ battery $\rightarrow$ load)
\end{itemize}

\section{Limitations and Assumptions}

\subsection{Model Assumptions}

\begin{enumerate}
    \item \textbf{Forecast Accuracy:} Model uses intraday forecasts (1-hour ahead) which have high accuracy
        \begin{itemize}
            \item \textbf{Approach:} Rolling horizon optimization refreshed every hour with updated ID forecasts
            \item \textbf{Reality:} ID forecast RMSE $\sim$70\% lower than DA forecasts
        \end{itemize}
    
    \item \textbf{Linear Efficiency:} Battery efficiency constant regardless of SOC or power level
        \begin{itemize}
            \item \textbf{Reality:} Efficiency varies with temperature, SOC, and C-rate
            \item \textbf{Impact:} $\pm$2--5\% energy throughput error
        \end{itemize}
    
    \item \textbf{No Battery Degradation:} Model doesn't account for capacity fade over time
        \begin{itemize}
            \item \textbf{Reality:} Capacity decreases $\sim$2\% per year (500--1000 cycles)
            \item \textbf{Mitigation:} Reduce $E^{\text{cap}}$ by 2\% annually in multi-year simulations
        \end{itemize}
    
    \item \textbf{No Network Constraints:} Assumes unlimited power transfer within REC
        \begin{itemize}
            \item \textbf{Reality:} Distribution lines have thermal limits
            \item \textbf{Mitigation:} Add line capacity constraints if critical
        \end{itemize}
    
    \item \textbf{Single Price Signal:} Uses day-ahead prices only
        \begin{itemize}
            \item \textbf{Reality:} Intraday prices, imbalance prices also relevant
            \item \textbf{Enhancement:} Multi-stage stochastic optimization
        \end{itemize}
\end{enumerate}

\section{References}

\subsection{Technical Standards}

\begin{enumerate}
    \item \textbf{IEC 61850:} Communication networks for substations (battery control protocols) \\
    \href{https://webstore.iec.ch/publication/6028}{https://webstore.iec.ch/publication/6028}
    
    \item \textbf{IEEE 1547:} Interconnection of distributed energy resources \\
    \href{https://standards.ieee.org/standard/1547-2018.html}{https://standards.ieee.org/standard/1547-2018.html}
    
    \item \textbf{VDE-AR-N 4110:} Technical requirements for DER grid connection (Germany) \\
    \href{https://www.vde.com/en/fnn/topics/technical-connection-rules}{https://www.vde.com/en/fnn/topics/technical-connection-rules}
\end{enumerate}

\subsection{Academic Literature}

\begin{enumerate}
    \item Peng, W., et al. (2018). ``Optimizing rooftop photovoltaic distributed generation with battery storage for peer-to-peer energy trading.'' \textit{Applied Energy}, 228, 2567--2580. \\
    DOI: \href{https://doi.org/10.1016/j.apenergy.2018.07.012}{10.1016/j.apenergy.2018.07.012}
    
    \item Castillo, A., \& Gayme, D. F. (2014). ``Grid-scale energy storage applications in renewable energy integration: A survey.'' \textit{Energy Conversion and Management}, 87, 885--894. \\
    DOI: \href{https://doi.org/10.1016/j.enconman.2014.07.063}{10.1016/j.enconman.2014.07.063}
    
    \item BDEW (Bundesverband der Energie- und Wasserwirtschaft). Standard load profiles documentation (Germany). \\
    \href{https://www.bdew.de/energie/standardlastprofile-strom/}{https://www.bdew.de/energie/standardlastprofile-strom/}
\end{enumerate}

\subsection{Software Tools}

\begin{enumerate}
    \item \textbf{Pyomo:} Python optimization modeling language \\
    \href{http://www.pyomo.org/}{http://www.pyomo.org/} | Documentation: \href{https://pyomo.readthedocs.io/}{https://pyomo.readthedocs.io/}
    
    \item \textbf{GLPK:} GNU Linear Programming Kit (open-source solver) \\
    \href{https://www.gnu.org/software/glpk/}{https://www.gnu.org/software/glpk/}
    
    \item \textbf{SimBench:} Benchmark datasets for power system analysis \\
    \href{https://simbench.de/en/}{https://simbench.de/en/} | Code: \href{https://github.com/e2nIEE/simbench}{https://github.com/e2nIEE/simbench}
\end{enumerate}

\section*{Document Control}

\begin{table}[h]
\centering
\begin{tabular}{@{}llll@{}}
\toprule
\textbf{Version} & \textbf{Date} & \textbf{Author} & \textbf{Changes} \\ 
\midrule
1.0 & 2026-02-01 & REC Optimization Team & Initial release \\
\bottomrule
\end{tabular}
\end{table}

\textbf{Approval:}
\begin{itemize}
    \item Technical Review: $\checkmark$ Pending
    \item Economic Analysis: $\checkmark$ Pending
    \item Regulatory Compliance: $\checkmark$ Pending
\end{itemize}

\textbf{Next Review Date:} 2026-05-01

\vspace{1cm}
\begin{center}
\textbf{--- END OF DOCUMENT ---}
\end{center}

\end{document}
